\begin{table}[htbp]
\centering
\caption{Minimum drug cost of treating everyone with HIV and keeping those at risk on PrEP, 2027--2046}
\label{tab:cost_to_end_hiv}
\resizebox{\textwidth}{!}{%
\begin{tabular}{lcccccc}
\toprule
 & \multicolumn{3}{c}{Present value, 3\%} & PV, 0\% & PV, 5\% & First-year \\
\cmidrule(lr){2-4}
Path: scenario & Treatment & PrEP & Total & Total & Total & cost \\
\midrule
Static: Lenacapavir only & 28.3 & 12.3 & 40.5 & 52.9 & 34.6 & 2.6 \\
Static: Lenacapavir + alimatravir & 28.3 & 7.3 & 35.6 & 46.2 & 30.5 & 2.6 \\
Dynamic: Lenacapavir only & 27.6 & 12.3 & 39.8 & 51.7 & 34.1 & 2.7 \\
Dynamic: Lenacapavir + alimatravir & 27.6 & 7.3 & 34.9 & 45.0 & 30.0 & 2.7 \\
\bottomrule
\end{tabular}}
\par\vspace{4pt}
\begin{minipage}{\textwidth}\footnotesize US\$ billion, drug prices only (current dollars, no inflation). Treatment: everyone living with HIV on generic TLD at \$45 per person-year. PrEP: 20 million people (UNAIDS) on generic lenacapavir at \$40; in the second scenario 50\% switch to alimatravir from 2029 at \$3, the estimated cost of production (Hill et al., AIDS 2026); Merck has not announced a price. Static: 2025 numbers held constant (41.0 million people with HIV). Dynamic: new infections fall linearly from 1.2 million (2025) to the 200,000 Fast-Track target in 2030, and the 2025 AIDS death rate among people with HIV is held constant; infections fall below deaths in 2029. Delivery costs (testing, visits, supply chain) are excluded, so these are lower bounds.\end{minipage}
\end{table}
